\documentclass[10pt]{article}

\usepackage[a4paper,margin=22mm]{geometry}
\usepackage[T1]{fontenc}
\usepackage{lmodern}
\usepackage{microtype}
\usepackage{graphicx}
\usepackage{booktabs}
\usepackage{xcolor}
\usepackage{titlesec}
\usepackage{caption}
\usepackage{enumitem}
\usepackage{array}
\usepackage[hidelinks]{hyperref}

\definecolor{secblue}{HTML}{1F4E8C}
\definecolor{rulegrey}{HTML}{888888}
\definecolor{codebg}{HTML}{F1F5F9}

\titleformat{\section}{\large\bfseries\color{secblue}}{\thesection}{0.6em}{}
\titleformat{\subsection}{\normalsize\bfseries\color{secblue}}{\thesubsection}{0.6em}{}
\titlespacing*{\section}{0pt}{12pt}{5pt}
\titlespacing*{\subsection}{0pt}{8pt}{3pt}

\captionsetup{font=small,labelfont=bf,labelsep=period}
\setlist[itemize]{leftmargin=1.4em,itemsep=2pt,topsep=3pt}
\renewcommand{\arraystretch}{1.15}

% inline code
\newcommand{\code}[1]{\texttt{\small #1}}

\setlength{\parindent}{0pt}
\setlength{\parskip}{5pt}

\begin{document}

% ---------------------------------------------------------------- title block
\begin{center}
  {\LARGE\bfseries cMCP --- Engineering Report}\\[4pt]
  {\normalsize A from-scratch Model Context Protocol implementation in pure C11,\\
   measured for spec conformance, performance, and hardening}\\[6pt]
  {\small\color{rulegrey} Generated 2026-06-09 \quad\textbullet\quad
   pure C11, three static link targets, no third-party JSON / MCP SDK}
\end{center}
\vspace{-2pt}
{\color{secblue}\hrule height 1.2pt}
\vspace{6pt}

% ================================================================= 1
\section{Executive summary}

cMCP is a from-scratch implementation of the Model Context Protocol (MCP) in
pure C11. It is a \emph{library}, not an application: it ships three static
link targets (\code{libcmcp\_core.a}, \code{libcmcp\_server.a},
\code{libcmcp\_client.a}) plus reference binaries. The JSON parser, JSON-RPC
framing, JSON-Schema subset validator, and all three transports are
hand-rolled; the only third-party dependency is \code{libcurl}, linked by the
HTTP \emph{client} alone. This report states what the implementation actually
does on the wire, how fast it is against the official reference SDKs, and how
it behaves under an adversarial peer --- and is careful about \emph{what} each
claim licenses.

\begin{itemize}
  \item \textbf{It is spec-faithful, and that is measured, not asserted.}
    cMCP is cross-checked against the TypeScript reference SDK in
    \emph{both} wire roles --- cMCP-client driving a TS server (32 assertions)
    and a TS client driving a cMCP server (8 checks) --- plus a captured-wire
    replay gate and a weekly spec-drift watch against the upstream spec.
  \item \textbf{It is fast where it matters and tiny at rest.} On the one
    apples-to-apples workload (\code{tools/call echo} over stdio, same client,
    server swapped) cMCP sustains \textbf{43\,494 calls/s} --- \textbf{5.4$\times$}
    the TypeScript SDK and \textbf{45$\times$} the Python SDK --- at a p50 of
    \textbf{21\,\textmu s}. A stdio server idles at \textbf{2.2\,MB} RSS, $27$--$33\times$
    lighter than the interpreted SDKs, from a 328\,KB static binary.
  \item \textbf{It is hardened against a hostile peer.} A STRIDE threat model
    drives concrete controls: an SSRF egress guard (DNS-rebinding-safe),
    schema validation \emph{before} dispatch, parse depth/element caps, stdio
    and HTTP frame caps, slowloris timeouts plus an accept-rate token bucket,
    and a credential log-redactor. Each maps to a named test.
  \item \textbf{The full quality gate passes clean.} 3\,885 assertions across
    30 test binaries, green under Valgrind, ASan/UBSan, and ThreadSanitizer;
    four libFuzzer harnesses; and a soak driver that shows no drift over long
    stdio and HTTP runs. Every gate is wired into CI.
  \item \textbf{It is scoped for stdio-primary, single-tenant embedding.}
    The design point is a host process (\textbf{butlerbot}) linking the client
    library to drive tools over stdio. HTTP is supported and hardened, but the
    network surface is \emph{defensive depth}; TLS and application-layer
    authentication are deliberately the reverse proxy's job (\S6), not a
    ``later'' item.
\end{itemize}

% ================================================================= 2
\section{Conformance: spec-faithful on the wire}

For a protocol library the central quality question is not ``does it run'' but
``does it produce \emph{exactly} the bytes the spec requires, and interoperate
with an independent implementation.'' cMCP answers this with cross-implementation
checks in both directions rather than self-consistency alone.

\begin{itemize}
  \item \textbf{cMCP client vs.\ TS reference server} (\code{make conformance},
    32 assertions): the cMCP async client runs the full handshake against the
    official TypeScript SDK server and exercises \code{tools/list},
    \code{tools/call}, \code{resources}, and \code{prompts}.
  \item \textbf{TS client vs.\ cMCP server} (8 checks): the reference TS client
    drives the cMCP \code{echo-server}, confirming the cMCP \emph{server} side
    is wire-compatible from the other direction.
  \item \textbf{Replay regression gate} (\code{make replay}): real wire
    transcripts captured through the \code{cmcp-tee} proxy are replayed and
    diffed, so a framing change cannot regress silently.
  \item \textbf{Spec-drift watch} (\code{make check-spec-drift}, weekly CI):
    compares the pinned \code{CMCP\_PROTOCOL\_VERSION} against the upstream
    spec directories and flags a divergence before it becomes a bug.
\end{itemize}

Two conformance details are worth calling out because they are easy to get
wrong and are pinned by test:

\textbf{Dual error channels.} A \code{tools/call} whose arguments fail schema
validation surfaces as a \emph{tool-level} result
(\code{\{isError:true, content:[\dots]\}}), not a JSON-RPC error --- matching
the TS reference SDK and wire-compatible with it. The \code{-32602} error
\emph{channel} instead carries protocol-level rejections (unknown tool,
malformed \code{protocolVersion}). A correct host must check both channels;
cMCP pins this with \code{test\_schema} and the host-probe.

\textbf{Version-tolerant handshake.} cMCP pins one protocol version per release
but does not abort when a peer advertises a different one: per spec the server
responds with success and advertises its own version, and the client captures
the server's version and lets the host decide whether to disconnect. Only a
missing or malformed \code{protocolVersion} yields a hard error.

\textbf{Schema is a documented subset.} The validator covers \code{type}
(incl.\ array-of-types), \code{properties}, \code{required},
\code{additionalProperties:false}, \code{enum}, \code{minLength}/\code{maxLength}
(Unicode code points), \code{minimum}/\code{maximum}, and \code{items}. It does
\emph{not} implement \code{\$ref}, \code{oneOf}/\code{anyOf}/\code{allOf},
\code{format}, or \code{pattern} in v0.1 --- a deliberate, documented scope
(\code{docs/schema-conformance.md}), not a silent gap.

% ================================================================= 3
\section{Performance vs.\ the state of the art}

The comparison below is deliberately apples-to-apples: the \emph{same} cMCP
client issues the \emph{same} \code{tools/call echo} workload over stdio, and
only the subprocess server is swapped per row. All three implementations pay
exactly one subprocess hop, so the measurement isolates the server runtime.
Numbers are 10\,000 iterations after a 1\,000-iteration warmup on the
development machine (\code{make bench-compare}, 2026-06-08).

\begin{figure}[ht]
  \centering
  \includegraphics[width=0.92\linewidth]{fig_throughput.png}
  \caption{Throughput on the one apples-to-apples engine comparison. Green is
  this work; grey are the official reference SDKs running the identical
  client and workload.}
\end{figure}

\begin{table}[ht]
  \centering
  \caption{Throughput, latency, and idle footprint. Latency is per round-trip;
  RSS is an idle stdio server.}
  \begin{tabular}{l r r r r r}
    \toprule
    Server & calls/s & vs cMCP & p50 (\textmu s) & p99 (\textmu s) & idle RSS \\
    \midrule
    \textbf{cMCP} (C11)      & \textbf{43\,494} & \textbf{1.0$\times$} & \textbf{21} & \textbf{31}   & \textbf{2.2\,MB} \\
    TypeScript SDK (Node)    & 8\,035           & 0.18$\times$         & 95          & 189           & 74\,MB \\
    Python SDK (FastMCP)     & 974              & 0.02$\times$         & 994         & 1\,315        & 60\,MB \\
    \bottomrule
  \end{tabular}
\end{table}

\begin{figure}[ht]
  \centering
  \includegraphics[width=0.78\linewidth]{fig_latency.png}\\[2pt]
  \includegraphics[width=0.78\linewidth]{fig_rss.png}
  \caption{Per-call latency (log axis) and idle memory footprint. cMCP's p99
  (31\,\textmu s) sits below the other implementations' p50; its idle RSS is a
  rounding error next to an interpreter and its \code{node\_modules}/venv.}
\end{figure}

\subsection{Why the gap --- and what the comparison licenses}

The spread is a \emph{runtime} story, not a protocol one. cMCP is native C11
with no interpreter, no GC, and (for a stdio server) zero third-party
dependencies: the server artifact is a 328\,KB static binary. The TypeScript
SDK carries a Node runtime plus \code{node\_modules} (\code{sdk}, \code{zod});
the Python SDK carries CPython plus a venv. The honest reading is that this is
an \textbf{engine / footprint comparison}, exactly the regime where a compiled
language wins, and \emph{not} a claim of feature parity --- the reference SDKs
ship a richer ecosystem (decorators, type-driven schema, framework
integrations) that cMCP intentionally does not. What the chart does license is
the deployment claim: at the \emph{butlerbot} embedding point, the protocol
layer costs microseconds and a couple of megabytes, so it disappears inside an
LLM inference cycle.

For reference, cMCP's pure in-process ceiling (no subprocess, inline pipe pair)
is \textbf{50\,487 calls/s at p50 = 19\,\textmu s}; the cross-SDK rows each pay
the subprocess hop so all three are measured identically. Because hardware and
kernel move these numbers, a weekly \code{bench-compare} CI job re-runs the
comparison and posts a fresh table to its run summary, so the published
snapshot never silently goes stale.

% ================================================================= 4
\section{Test posture and quality gates}

cMCP's suites map onto the standard Unit $\rightarrow$ Integration $\rightarrow$
Evaluation pyramid, with a security/hardening axis that cuts across all layers.
The hermetic \code{make test} run is 3\,885 assertions across 30 binaries.

\begin{figure}[ht]
  \centering
  \includegraphics[width=0.92\linewidth]{fig_pyramid.png}
  \caption{Where the assertions live. L3 (cross-implementation evaluation) is
  opt-in --- it needs Node and network --- so it is excluded from the hermetic
  count and reported separately in \S2.}
\end{figure}

\begin{table}[ht]
  \centering
  \caption{Quality gates, all wired into CI; green as of 2026-06-08.}
  \begin{tabular}{l l l}
    \toprule
    Gate & Command & What it proves \\
    \midrule
    Unit + integration   & \code{make test}            & 3\,885 assertions / 30 binaries \\
    Memory safety        & \code{make valgrind}        & leak-free, no invalid access \\
    ASan + UBSan         & \code{make test-asan}       & no heap / undefined-behaviour errors \\
    ThreadSanitizer      & \code{make test-tsan}       & no data races \\
    Fuzzing              & \code{make fuzz-smoke}      & 4 libFuzzer harnesses (json/rpc/schema/http) \\
    Wire regression      & \code{make replay}          & captured-transcript replay diff \\
    Soak (stdio + HTTP)  & \code{make soak}            & no drift over long runs \\
    Cross-impl conformance & \code{make conformance}   & matches TS reference SDK, both roles \\
    Spec-drift watch     & \code{make check-spec-drift}& pinned protocol vs.\ upstream spec \\
    \bottomrule
  \end{tabular}
\end{table}

The integration tier is not mocked: \code{test\_stdio\_roundtrip} actually
\code{fork()}s a child server and runs a real handshake, and
\code{test\_client\_server} wires real \code{cmcp\_server\_run} instances onto
pipe pairs to exercise the async client and multi-server session end-to-end.
The concurrency model under test is one reader thread per transport, handlers
on a small worker pool, and a single writer mutex per transport --- TSan-clean.

% ================================================================= 5
\section{Robustness and security (hardening)}

cMCP is a protocol library, so the trust model is precise: the host is trusted
(cMCP runs in its address space), but the \emph{peer} on the other end of a
transport may be hostile or buggy and its input is untrusted. The hardening
maps a STRIDE threat catalogue onto concrete, individually-tested controls.

\begin{table}[ht]
  \centering
  \caption{Threat $\rightarrow$ control $\rightarrow$ test. Every control is
  reachable from an environment knob and pinned by a named suite.}
  \begin{tabular}{p{0.30\linewidth} p{0.42\linewidth} p{0.20\linewidth}}
    \toprule
    Threat & Control & Test \\
    \midrule
    SSRF via discovery & egress guard on the \emph{resolved} peer (rejects
      metadata / link-local / RFC1918 / CGNAT / ULA); DNS-rebinding-safe
      & \code{test\_http\_ssrf\_guard} \\
    DNS rebinding & \code{Origin} allow-list on the HTTP server
      & \code{test\_http\_server} \\
    Insecure deserialization & schema validation \emph{before} dispatch; parse
      depth + per-container element caps
      & \code{test\_schema}, \code{test\_hostile\_peer} \\
    Memory-bomb / OOM & stdio frame cap, HTTP body cap, handler \code{RLIMIT\_AS}
      & \code{test\_stdio\_frame\_cap}, \code{test\_handler\_rlimit} \\
    Slowloris / spike & accept-rate token bucket, per-recv idle timeout +
      cumulative deadline
      & \code{test\_hardening\_cross\_knob} \\
    Credential leakage & log redactor over credential-shaped keys
      & \code{test\_http\_auth}, \code{test\_logging} \\
    Malformed / hostile peer & bounded parsers, fail-closed framing
      & \code{test\_hostile\_peer} \\
    \bottomrule
  \end{tabular}
\end{table}

The controls are layered front-to-back: a request crosses the \code{Origin}
allow-list and SSRF guard, then frame / body / depth / element caps, then
schema validation, then rate-limit and idle/deadline timeouts, before a handler
runs on the worker pool under \code{RLIMIT\_AS}; anything the handler logs
passes the credential redactor on the way out. Each is governed by a snapshot-once
environment knob (e.g.\ \code{CMCP\_HTTP\_ALLOW\_PRIVATE},
\code{CMCP\_JSON\_MAX\_DEPTH}, \code{CMCP\_STDIO\_MAX\_FRAME},
\code{CMCP\_HTTP\_ACCEPT\_RATE}, \code{CMCP\_LOG\_REDACT}), so an operator can
tune or disable a control for an environment they own without recompiling.
Beyond the unit gates, four libFuzzer harnesses drive the JSON, JSON-RPC,
schema, and HTTP-head parsers against curated corpora, and the soak driver runs
the stdio and HTTP servers for long windows (with optional server respawn) to
prove no leak or drift accumulates.

% ================================================================= 6
\section{Operating envelope --- scope and deliberate non-goals}

A report owes an honest answer to ``what is this \emph{not} for.'' cMCP's design
point is a host process linking the client library to drive one or more MCP
servers over stdio, single-tenant. The network surface exists and is hardened,
but two large areas are deliberately out of scope --- by architecture, not by
omission.

\textbf{TLS is terminator-only.} cMCP links no TLS library. A protocol library
is the wrong layer to own cert rotation, OCSP stapling, and session resumption;
the state-of-the-art shape for a JSON-RPC server is behind nginx / caddy /
Envoy, which already do TLS, HTTP/2, and rate-limiting better than a
single-purpose library could, and linking a TLS stack would force one
implementation choice on every downstream. This is an architectural decision
for cMCP's lifetime, not a ``TLS later'' position; the proxy recipes (including
\code{X-Forwarded-Proto} trust) live in \code{docs/deployment-tls.md}.

\textbf{Application-layer authentication is the terminator's job.} OAuth\,2.1 /
API keys / mTLS user identification happen at the proxy; cMCP receives the
already-authenticated request and trusts it. Likewise out of scope: handler-level
resource exhaustion (disk, egress, child processes --- the handler decides
what is safe), cryptographic correctness beyond the single CSPRNG read used to
mint session ids, and timing side channels (constant-time session-id compare is
not implemented, a documented low-risk gap that requires near-position
privilege to exploit).

cMCP is intentionally \emph{not} cRAG-shaped: cRAG is one example consumer
(wrapped by the in-tree \code{crag-mcp} reference server), not a privileged
dependency, and no cRAG-specific code path lives inside \code{src/}.

% ================================================================= 7
\section{Methodology and honest caveats}

\begin{itemize}
  \item \textbf{Performance numbers are machine-specific.} They come from
    \code{make bench-compare} on the development box; kernel, build flags, and
    hardware move them. The weekly \code{bench-compare} CI job re-runs the
    comparison so the snapshot stays current; treat absolute values as
    indicative and the \emph{ratios} as the durable signal.
  \item \textbf{The cross-SDK comparison is engine-level.} All three rows run
    the identical client and workload and pay one subprocess hop, so the
    comparison is fair --- but it measures runtime and footprint, not feature
    parity. The reference SDKs offer ecosystem conveniences cMCP does not.
  \item \textbf{L3 evaluation is opt-in.} The cross-implementation conformance
    and replay checks need Node and network, so they sit outside the hermetic
    \code{make test} gate and are reported separately. The 3\,885-assertion
    figure is the offline floor, not the ceiling.
  \item \textbf{The schema validator is a documented v0.1 subset.} Unsupported
    keywords (\code{\$ref}, \code{oneOf}/\code{anyOf}/\code{allOf},
    \code{format}, \code{pattern}) are listed, not silently dropped.
  \item \textbf{The threat model is stdio-primary.} HTTP-axis controls
    (SSRF, origin, slowloris, accept-rate) are defensive depth for a surface
    that, in the target deployment, sits behind a trusted terminator.
\end{itemize}

\vfill
{\color{rulegrey}\hrule height 0.6pt}
\vspace{4pt}
{\small\color{rulegrey}
cMCP --- pure C11, three static link targets, no third-party JSON / MCP SDK.
All figures reproducible: \code{make test}, \code{make bench-compare},
\code{make conformance}, \code{make replay}.}

\end{document}
